\documentclass[lean]{AFM}
\graphicspath{ {./images/} }

\addbibresource{krafft.bib}
\setcounter{biburllcpenalty}{7000}
\setcounter{biburlucpenalty}{8000}


\renewcommand{\.}{\hskip .75pt}
\renewcommand{\arraystretch}{1.25}
\newcommand{\fin}[1]{[\.#1\.]}
\newcommand{\sekt}[1]{Section~\ref{#1}}

\def\urlpfx{https://github.com/ElNando888/KrafftSieve/blob/v1.0}

\makeatletter
\newcommand{\thickhline}{%
	\noalign {\ifnum 0=`}\fi \hrule height 2pt
	\futurelet \reserved@a \@xhline
}
\newcolumntype{"}{@{\hskip\tabcolsep\vrule width 2pt\hskip\tabcolsep}}
\makeatother

\makeatletter
\newcommand{\customlabel}[2]{%
   \protected@write \@auxout {}{\string \newlabel {#1}{{#2}{\thepage}{#2}{#1}{}} }%
   \hypertarget{#1}{#2\!\!}
}
\makeatother


\title[a short title]{The title}

\author[F. Portela]{Fernando Portela} 

\authorinfo[F. Portela]{Independent Researcher}{math.nt@elnando.anonaddy.me}

\VOLUME{0}
\YEAR{2026}
\NUMBER{0}
\firstpage{1}
\DOI{https://doi.org/10.46298/afm.?????}
\receiveddate{00}
\finaldate{00}
\accepteddate{00}
\msc{11N05, 11N35, 11N36, 03B35}



\begin{abstract}

...

\end{abstract}

\keywords{Krafft form, Twin primes, Turán sieve, Interactive theorem proving, Lean 4}

\begin{document}


\section{Introduction} \label{introduction}

.

\section {first section} \label{first}

.

\section {Conclusion}

.

\section*{Acknowledgements}
This work would not have been possible without the tireless assistance of the Aristotle
framework and Gemini 3.1 Pro. The author wishes to express his gratitude for their
unwavering dedication to parsing the intricate complexities of the formalization,
illuminating structural insights, and maintaining a steadfast, collaborative dialogue. More
than mere computational instruments, they served as indispensable collaborators whose tireless
deductive logic helped chart the course through these formalized geometric spaces. Specifically,
Gemini 3.1 Pro was instrumental in clarifying the theoretical ideas and generating precise prompt
heuristics for the formal prover, while Aristotle served as the primary mechanical workhorse
securing the dense derivations.

\printbibliography

\newpage

\section*{Appendix: Dependencies between theorems}

The following directed graphs map the logical dependencies between the central formalized results.
To maintain structural clarity and avoid overwhelming edge density, basic definitions and
foundational types are generally excluded from the hierarchy. 

\begin{itemize}
    \setlength\itemsep{0em}
    \item \textbf{Blue-shaded rounded rectangles}: Formal proofs (Lemmas and Theorems).
    \item \textbf{Green-shaded rectangles}: Central definitions or structures.
    \item \textbf{Directed arrows} ($A \rightarrow B$): Logical dependence, signifying that proof
    $A$ explicitly ``uses'' or ``applies'' result $B$ to close a goal.
\end{itemize}

\begin{figure}[h]
\centering
\includegraphics[width=0.7\textwidth]{theorem1.png}
\caption{.}
\label{fig:thmtree1}
\end{figure}

\end{document}